\SourceLeaf{302}
\par\bigskip\pdfbookmark[0]{Bemerkungen zu De nexu}{nexus-bemerkungen}
\begin{center}
\textsc{Bemerkungen zur Abhandlung}\\[1em]
\textsc{De nexu inter multitudinem classium, in quas formae binariae secundi gradus distribuuntur, earumque determinantem.}
\end{center}
\BookRule{.16\textwidth}
\par\medskip\pdfbookmark[1]{Zu I. und II.}{nexus-note-I-II}\centerline{Zu I. und II.}\smallskip

Die zweite Formel für die Anzahl der innerhalb des Kreises liegenden Puncte (I. art. 3 und II. art. 5) ergiebt sich aus der Betrachtung des in denselben eingeschriebenen Quadrates, dessen Seiten den Coordinatenaxen parallel sind; die Vergleichung beider Formeln führt zu dem auch arithmetisch leicht zu beweisenden Satze
\[
r'+r''+\ldots+r^{(q)}=qq+r^{(q+1)}+r^{(q+2)}+\ldots+r^{(r)}
\]
aus welchem sich wieder die Richtigkeit der ersten von den beiden folgenden Regeln ergiebt, die sich auf einem besondern Blatt vorfanden:

„Auflösungen der Gleichung $xx+yy\leqq A$; formula
\[
1+4\surd A+4\surd\tfrac12 A+8\Sigma(\surd(A-nn)-n)
\]
wo bei jeder Wurzel der Bruch weggelassen und von $n=1$ bis $n=\surd A$ „(soll heissen $\surd\tfrac12 A$)“ summirt wird.

Andre Formel
\[
1+4\{A-\tfrac A3+\tfrac A5-\tfrac A7+\tfrac A9-\tfrac A{11}\ldots\}
\]
wo bei jedem Theil der Bruch weggelassen.“

Diese letztere Formel folgt aus dem später (I. art. 6) zur Anwendung kommenden Satze über die Anzahl aller verschiedenen Darstellungen einer bestimmten Zahl durch die Form $xx+yy$ (vergl. Disqq. Arithm. art. 182, Note), welcher leicht in den folgenden umgeformt werden kann: die Anzahl der verschiedenen Darstellungen einer positiven ganzen Zahl $m$ durch die Form $xx+yy$ ist $=4(a-b)$, wo $a,b$ die Anzahlen der Divisoren von $m$ bedeuten, welche resp. von der Form $4n+1,\ 4n+3$ sind. Aus der Vergleichung
\SourceLeaf{303}
dieser arithmetischen Formel mit der (in I. art. 5 oder II. art. 4) durch geometrische Betrachtungen gewonnenen mittlern Darstellungsanzahl erhält man leicht und in aller Strenge das bekannte Resultat
\[
\frac\pi4=1-\tfrac13+\tfrac15-\tfrac17+\ldots
\]
welches in der Abhandlung (I. art. 7) durch eine ähnliche Vergleichung, aber mit Hülfe unendlicher Producte abgeleitet wird.
\BookRule{.25\textwidth}
